%! Constructing a Confidence Interval for a Population Proportion — Unit 6, Lesson 6.2 — Homework (Answer Key)
\documentclass[10pt]{article}
\usepackage{apstats-article}
\usepackage{apstats-key}   % pulls in -boxes; do NOT also load -boxes
\usepackage{grading}
\gradingassignment{apstats-u06-l02-hw}
\gradingcourse{advanced_statistics}
\newcommand{\TallMath}[1]{$\displaystyle #1\rule[-1.4em]{0pt}{3.2em}$}

\begin{document}

\pageheader{Unit 6, Lesson 6.2}{Homework --- Answer Key}
{\small\itshape Write each final answer in its box. Only what is inside a box is graded.}\par

\vspace{0.15in}

\textit{For each problem, use the four-step procedure: State, Plan (verify conditions), Do (compute interval), Conclude (interpret).}

\begin{enumerate}

  \item A random sample of 400 registered voters in a city finds that 228 support a proposed
        park expansion. Construct a 95\% confidence interval for the true proportion of all
        registered voters in the city who support the expansion.

        \textbf{State:} \gradedhide{\writeline\\\gradedhere{\gradedblank[id=hw1a, type=rubric, points=1, display={\bfseries $p$ = true proportion of all registered voters in the city who support the expansion}]{\linewidth+6pt}{\gradedlinesheight{9}}}\writelines{10}}

        \textbf{Plan:}
        \begin{itemize}
          \item Random: \gradedhide{\writeline\\\gradedhere{\gradedblank[id=hw1b1, type=rubric, points=1, display={\bfseries Random sample (stated): met}]{\linewidth+6pt}{\gradedlinesheight{9}}}\writelines{9}}
          \item Large Counts: \gradedhide{\writeline\\\gradedhere{\gradedblank[id=hw1b2, type=rubric, points=1, display={\bfseries $228 \ge 10$ and $172 \ge 10$: met}]{\linewidth+6pt}{\gradedlinesheight{9}}}\writelines{9}}
          \newpage
          \item 10\% rule: \gradedhide{\writeline\\\gradedhere{\gradedblank[id=hw1b3, type=rubric, points=1, display={\bfseries $400 \le 10\%$ of all voters: met}]{\linewidth+6pt}{\gradedlinesheight{9}}}\writelines{9}}
        \end{itemize}

        \textbf{Do:} $\hat p = $ \smash{\gradedblank[id=hw1c1, type=short, points=1, answer=0.57, check=numeric]{2cm}{0.17in}}
        \[
          CI = \ans{0.57} \pm \smash{\gradedblank[id=hw1c2, type=short, points=1, answer=1.960, check=numeric, tolerance=0.005]{1.5cm}{0.17in}}\smash{\sqrt{\frac{\ans{0.57(0.43)}}{\ans{400}}}}\vphantom{\sqrt{\frac{\blank{2cm}}{\blank{1.5cm}}}}
             = (\smash{\gradedblank[id=hw1d1, type=short, points=1, answer=0.521, check=numeric, tolerance=0.002]{2.5cm}{0.17in}},\; \smash{\gradedblank[id=hw1d2, type=short, points=1, answer=0.619, check=numeric, tolerance=0.002]{2.5cm}{0.17in}})
        \]

        \textbf{Conclude:} \gradedhide{\writeline\\\gradedhere{\gradedblank[id=hw1e, type=rubric, points=2, display={95\% confident (0.521, 0.619) captures the true proportion of city voters who support the expansion.}]{\linewidth+6pt}{\gradedlinesheight{1}}}\writelines{9}}

  \newpage
  \item A wildlife biologist tags and releases 600 randomly selected fish from a lake.
        Later, a random sample of 80 fish from the lake is taken, and 12 are tagged.
        Construct a 99\% confidence interval for the true proportion of tagged fish in the lake.

        \textbf{State:} \gradedhide{\writeline\\\gradedhere{\gradedblank[id=hw2a, type=rubric, points=1, display={\bfseries $p$ = true proportion of fish in the lake that are tagged}]{\linewidth+6pt}{\gradedlinesheight{9}}}\writelines{10}}

        \textbf{Plan:}
        \begin{itemize}
          \item Random: \gradedhide{\writeline\\\gradedhere{\gradedblank[id=hw2b1, type=rubric, points=1, display={\bfseries Random sample (stated): met}]{\linewidth+6pt}{\gradedlinesheight{9}}}\writelines{9}}
          \item Large Counts: \gradedhide{\writeline\\\gradedhere{\gradedblank[id=hw2b2, type=rubric, points=1, display={\bfseries $12 \ge 10$ and $68 \ge 10$: met}]{\linewidth+6pt}{\gradedlinesheight{9}}}\writelines{9}}
          \newpage
          \item 10\% rule: \gradedhide{\writeline\\\gradedhere{\gradedblank[id=hw2b3, type=rubric, points=1, display={\bfseries $80 \le 10\%$ of all fish: met}]{\linewidth+6pt}{\gradedlinesheight{9}}}\writelines{9}}
        \end{itemize}

        \textbf{Do:}
        \[
          CI = (\smash{\gradedblank[id=hw2c1, type=short, points=1, answer=0.047, check=numeric, tolerance=0.002]{2.5cm}{0.17in}},\; \smash{\gradedblank[id=hw2c2, type=short, points=1, answer=0.253, check=numeric, tolerance=0.002]{2.5cm}{0.17in}})
        \]

        \textbf{Conclude:} \gradedhide{\writeline\\\gradedhere{\gradedblank[id=hw2d, type=rubric, points=2, display={99\% confident (0.047, 0.253) captures the true proportion of tagged fish in the lake.}]{\linewidth+6pt}{\gradedlinesheight{1}}}\writelines{9}}

  \newpage
  \item A school district wants to estimate the proportion of families who would use
        an after-school tutoring program. An administrator wants a 95\% confidence interval
        with a margin of error of at most 0.05. A pilot survey suggests $\hat p \approx 0.40$.
        \begin{enumerate}[label=\alph*.]
          \item Find the minimum required sample size using the pilot estimate.
                \[n \ge \left(\frac{z^*}{m}\right)^2 \hat p(1-\hat p) = \smash{\gradedblank[id=hw3a1, type=short, points=1, answer=(1.96/0.05)^2(0.40)(0.60), check=expression, display={$1536.64 \times 0.24$}]{6cm}{0.17in}} = \smash{\gradedblank[id=hw3a2, type=short, points=2, answer=368.79, check=numeric, tolerance=0.5]{2.5cm}{0.17in}}\]
                Minimum $n = $ \smash{\gradedblank[id=hw3a3, type=short, points=1, answer=369, check=numeric, tolerance=0]{2cm}{0.17in}}

          \item How would your answer change if you used $\hat p = 0.50$ instead? Would it be
                larger or smaller? \gradedhide{\writeline\\\gradedhere{\gradedblank[id=hw3b, type=rubric, points=2, display={Larger: $1536.64(0.25) = 384.16$, so $n = 385$; $\hat p(1-\hat p)$ is largest at 0.5.}]{\linewidth+6pt}{\gradedlinesheight{1}}}\writelines{9}}
        \end{enumerate}

\end{enumerate}

\begin{reflectionbox}
\textbf{Preview:} In Lesson 6.3 you will focus on \emph{interpreting} confidence intervals ---
what ``95\% confident'' really means for the method, and how to use an interval to evaluate a
researcher's claim about a population proportion. Think: if your CI does \emph{not} contain a
claimed value, what conclusion can you draw?
\end{reflectionbox}

\end{document}
